% !TeX root = ../main.tex

\chapter{从反常霍尔效应到非线性反常霍尔效应}

\section{线性反常霍尔电导率}
\label{sec:AHE}

反常霍尔效应（AHE）指铁磁材料中无外加磁场时出现的横向霍尔电压。电流$\mathbf{J}$与电场$\mathbf{E}$由电导率张量联系：
\begin{equation}
J_i = \sigma_{ij} E_j,
\end{equation}
其中霍尔响应来自反对称部分$\sigma_{ij}^A = \frac{1}{2}(\sigma_{ij}-\sigma_{ji})$，三维下可用轴矢量表示为$\sigma_{ij}^A = \epsilon_{ijk}\sigma_k^H$。

内禀贡献源于动量空间的贝里曲率。在自旋-轨道耦合和时间反演对称性破缺的体系中，约定$e>0$为元电荷，电子电荷为$-e$，运动方程为$\hbar\dot{\mathbf{k}}=-e\mathbf{E}$。布洛赫电子的反常速度为
\begin{equation}
\mathbf{v}^{\text{an}} = \frac{e}{\hbar}\mathbf{E}\times\bm{\Omega}_n(\mathbf{k}),
\end{equation}
由此直接给出内禀反常霍尔电导率（完整半经典推导见附录~\ref{nonlinear-hall}）：
\begin{equation}
\sigma_{ij} = -\frac{e^2}{\hbar}\epsilon_{ijk} M_k, \qquad
M_k = \sum_n \int \frac{d^3k}{(2\pi)^3} f_n^{(0)}(\mathbf{k})\,\Omega_{n,k}(\mathbf{k}).
\label{eq:AHE_sigma}
\end{equation}
贝里曲率的定义与量子几何关系详见\autoref{berry}。

\section{线性反常霍尔效应的对称性分析}
\label{sec:AM_AHE_symmetry}

共线/共面磁体中，线性AHE的存在性由单一判据决定：计入SOC后体系是否属于自旋-轨道磁体（SOM）。该判据与是否交变磁体无关。

\begin{itemize}
    \item \textbf{无SOC极限（SSG）}：所有共线或共面磁体的SSG同时强制自旋磁化和轨道磁化为零，线性AHE被严格禁戒。
    \item \textbf{计入SOC后（MSG）}：对称性降为磁空间群。若MSG不强制净磁化为零，则体系为SOM，线性AHE被允许且与轨道磁化同阶涌现；若MSG仍强制磁化为零，则为纯反铁磁体，线性禁止。
    \item \textbf{与交变磁性的关系}：交错磁体的本质是非相对论自旋劈裂；而反常霍尔效应来自于引入SOC后的对称性降低。
\end{itemize}
\vspace{5mm}
\textbf{小结：}
\begin{center}
\begin{tabular}{c c c}
\toprule
& 无SOC（SSG） & 有SOC（MSG） \\
\midrule
纯AFM & 禁止 & 禁止 \\
SOM   & 禁止 & \textbf{AHE允许} \\
\bottomrule
\end{tabular}
\end{center}
\vspace{5mm}
表~\ref{tab:3d} 和表~\ref{tab:2d} 列出了三维和二维磁点群中由贝里曲率多极子诱导的主导本征反常霍尔响应阶次。一阶（线性）AHE对应的磁点群恰好是SOM中MSG不强制磁化为零的铁磁类点群；禁戒一切本征霍尔响应的为包含$PT$对称或强制磁化为零的纯反铁磁类点群；其余磁点群允许高阶非线性霍尔效应。

\begin{table}[H]
\centering
\caption{所有122个三维磁点群中由Berry曲率多极子诱导的主导本征反常霍尔响应}
\label{tab:3d}
\begin{tabular}{p{3.5cm} p{10cm}}
\toprule
\textbf{主导响应阶次} & \textbf{磁点群} \\
\midrule
无反常霍尔响应 & 
$\bar{1}1'$, $\bar{1}'$, $2/m1'$, $2'/m$, $2/m'$, $mmm1'$, $m'mm$, $m'm'm'$, $4/m1'$, $4/m'$, $4'/m'$, $4/mmm1'$, $4/m'mm$, $4'/m'm'm$, $4/m'm'm'$, $\bar{3}1'$, $\bar{3}'$, $\bar{3}m1'$, $\bar{3}'m$, $\bar{3}'m'$, $6/m1'$, $6'/m$, $6/m'$, $6/mmm1'$, $6/m'mm$, $6'/mmm'$, $6/m'm'm'$, $m\bar{3}1'$, $m'\bar{3}'$, $m\bar{3}m1'$, $m'\bar{3}'m$, $m'\bar{3}'m'$ \\
\midrule
一阶（线性反常霍尔） & 
$1$, $\bar{1}$, $2$, $2'$, $m$, $m'$, $2/m$, $2'/m'$, $2'2'2$, $m'm2'$, $m'm'2$, $m'm'm$, $4$, $\bar{4}$, $4/m$, $42'2'$, $4m'm'$, $\bar{4}2'm'$, $4/mm'm'$, $3$, $\bar{3}$, $32'$, $3m'$, $\bar{3}m'$, $6$, $\bar{6}$, $6/m$, $62'2'$, $6m'm'$, $\bar{6}m'2'$, $6/mm'm'$ \\
\midrule
二阶（非线性霍尔） & 
$11'$, $21'$, $m1'$, $222$, $2221'$, $mm2$, $mm21'$, $41'$, $4'$, $\bar{4}1'$, $\bar{4}'$, $422$, $4221'$, $4'22'$, $4mm$, $4mm1'$, $4'm'm$, $\bar{4}2m$, $\bar{4}2m1'$, $\bar{4}2'm$, $\bar{4}2m'$, $31'$, $32$, $321'$, $3m$, $3m1'$, $61'$, $6'$, $622$, $6221'$, $6'22'$, $6mm$, $6mm1'$, $6'mm'$, $23$, $231'$, $432$, $4321'$, $4'32'$ \\
\midrule
三阶（非线性霍尔） & 
$mmm$, $4'/m$, $4/mmm$, $4'/mmm'$, $\bar{3}m$, $\bar{6}'$, $6'/m'$, $\bar{6}m2$, $\bar{6}'m'2$, $\bar{6}'m2'$, $6/mmm$, $6'/m'mm'$, $m\bar{3}$, $\bar{4}'3m'$, $m\bar{3}m'$ \\
\midrule
四阶（非线性霍尔） & 
$\bar{6}1'$, $\bar{6}m21'$, $\bar{4}3m$, $\bar{4}3m1'$ \\
\midrule
五阶（非线性霍尔） & 
$m\bar{3}m$ \\
\bottomrule
\end{tabular}
\end{table}

\begin{table}[H]
\centering
\caption{所有31个二维磁点群中由Berry曲率多极子诱导的主导本征反常霍尔响应}
\label{tab:2d}
\begin{tabular}{p{3.5cm} p{8cm}}
\toprule
\textbf{主导响应阶次} & \textbf{磁点群} \\
\midrule
无反常霍尔响应 & 
$21'$, $2'$, $2mm1'$, $2'mm'$, $41'$, $4mm1'$, $61'$, $6'$, $6mm1'$, $6'm'm$ \\
\midrule
一阶（线性反常霍尔） & 
$1$, $m'$, $2$, $2m'm'$, $4$, $4m'm'$, $3$, $3m'$, $6$, $6m'm'$ \\
\midrule
二阶（非线性霍尔） & 
$11'$, $m$, $m1'$ \\
\midrule
三阶（非线性霍尔） & 
$2mm$, $4'$, $4'm'm$ \\
\midrule
四阶（非线性霍尔） & 
$3m$, $31'$, $3m1'$ \\
\midrule
五阶（非线性霍尔） & 
$4mm$ \\
\midrule
七阶（非线性霍尔） & 
$6mm$ \\
\bottomrule
\end{tabular}
\end{table}

\clearpage

\section{非线性电导率的一般公式}
\label{sec:nonlinear_formula}

非线性霍尔效应自提出以来在理论和实验方面均取得了重要进展
\cite{Sodemann2015,gao_field_2014,nandy_symmetry_2019,du_quantum_2021,du_nonlinear_2021,ma_observation_2019,kang_nonlinear_2019,wang_intrinsic_2021,ba_nonlinear_2023,min_colossal_2024,lai_third-order_2021,Zhang2023,gao_quantum_2023,yu_quantum_2025,liu_giant_2025,li_quantum_2024,Fang2024,chu_third-order_2025,chen_giant_2026,zhu_magnetic_2025}。
当线性AHE被对称性禁戒时，高阶非线性霍尔效应成为主导。本节直接给出二阶与三阶电导率的核心公式，完整推导见附录~\ref{nonlinear-hall}。

\subsection{符号约定}
$a,b,c,d\in\{x,y,z\}$为坐标索引； $\partial_{k^a}\equiv \partial/\partial k_a$； $\epsilon_n$为能带色散； $f_n$为费米-狄拉克分布； $\Omega_n^{ab}$为贝里曲率； $G_n^{ab}$为量子度规； $\sum_n\int_{\mathbf{k}}$表示对占据能带求和并积分布里渊区。

\subsection{二阶电导率}
$j_c^{(2)} = \sigma^{ab;c}E_a E_b$，$\sigma^{ab;c}$对$(a,b)$对称。
\begin{align}
\sigma_{\text{Drude}}^{ab;c} &= -\frac{e^3\tau^2}{\hbar^3}\sum_n\int_{\mathbf{k}} f_n\,\partial_{k^a}\partial_{k^b}\partial_{k^c}\epsilon_n, \label{eq:sigma2_Drude}\\
\sigma_{\text{BCD}}^{ab;c} &= \frac{e^3\tau}{\hbar^2}\sum_n\int_{\mathbf{k}} f_n\,
\frac{1}{2}\bigl(\partial_{k^a}\Omega_n^{bc}+\partial_{k^b}\Omega_n^{ac}\bigr), \label{eq:sigma2_BCD}\\
\sigma_{\text{QMD}}^{ab;c} &= \frac{e^3}{\hbar}\sum_n\int_{\mathbf{k}} f_n\,
\Bigl(2\partial_{k^c}G_n^{ab}-\frac{1}{2}(\partial_{k^a}G_n^{bc}+\partial_{k^b}G_n^{ac})\Bigr). \label{eq:sigma2_QMD}
\end{align}

\subsection{三阶电导率}
$j_d^{(3)} = \sigma^{abc;d}E_a E_b E_c$，$\sigma^{abc;d}$对$(a,b,c)$全对称。
\begin{align}
\sigma_{\text{Drude}}^{abc;d} &= \frac{e^4\tau^3}{\hbar^4}\sum_n\int_{\mathbf{k}} f_n\,\partial_{k^a}\partial_{k^b}\partial_{k^c}\partial_{k^d}\epsilon_n, \label{eq:sigma3_Drude}\\
\sigma_{\text{BCQ}}^{abc;d} &= -\frac{e^4\tau^2}{\hbar^3}\sum_n\int_{\mathbf{k}} f_n\,
\frac{1}{3}\bigl(\partial_{k^a}\partial_{k^b}\Omega_n^{cd} + \partial_{k^b}\partial_{k^c}\Omega_n^{ad} + \partial_{k^a}\partial_{k^c}\Omega_n^{bd}\bigr), \label{eq:sigma3_BCQ}\\
\sigma_{\text{QMQ}}^{abc;d} &= \frac{e^4\tau}{\hbar^2}\sum_n\int_{\mathbf{k}} f_n\,
\frac{1}{3}\Bigl[\,2\bigl(\partial_{k^a}\partial_{k^d}G_n^{bc}+\partial_{k^b}\partial_{k^d}G_n^{ac}+\partial_{k^c}\partial_{k^d}G_n^{ab}\bigr) \nonumber\\
&\qquad -\bigl(\partial_{k^a}\partial_{k^c}G_n^{bd}+\partial_{k^b}\partial_{k^c}G_n^{ad}+\partial_{k^a}\partial_{k^b}G_n^{cd}\bigr)\Bigr], \label{eq:sigma3_QMQ}\\
\sigma_{\text{AIC}}^{abc;d} &= \frac{e^4}{\hbar}\sum_n\int_{\mathbf{k}} f_n\,
\frac{2}{3\epsilon_{n\bar{n}}}\bigl(G_n^{ab}\Omega_n^{cd}+G_n^{ac}\Omega_n^{bd}+G_n^{bc}\Omega_n^{ad}\bigr). \label{eq:sigma3_AIC}
\end{align}

\subsection{$\tau$ 标度与物理来源}
\begin{table}[H]
\centering
\caption{非线性电导率各贡献项的$\tau$标度与物理来源}
\label{tab:tau_scaling}
\begin{tabular}{c|c|c|c}
\toprule
响应阶数 & 贡献项 & $\tau$依赖 & 物理来源 \\
\midrule
\multirow{3}{*}{二阶} & Drude & $\tau^2$ & 能带色散 \\
                      & 贝里曲率偶极矩 & $\tau^1$ & 贝里曲率梯度 \\
                      & 量子度规偶极矩 & $\tau^0$ & 量子度规梯度 \\
\midrule
\multirow{4}{*}{三阶} & Drude & $\tau^3$ & 能带色散 \\
                      & 贝里曲率四极矩 & $\tau^2$ & 贝里曲率二阶梯度 \\
                      & 量子度规四极矩 & $\tau^1$ & 量子度规二阶梯度 \\
                      & 额外带间贡献 & $\tau^0$ & 度规-曲率乘积 \\
\bottomrule
\end{tabular}
\end{table}
\section{量子几何量在对称操作下的变换规则}

本节用投影算符语言推导动量导数 \(\partial_a\)、带归一化量子度规 \(G^{ab}\)、
贝里曲率张量 \(\Omega^{ab}\) 及其赝矢量对偶 \(\Omega^c\) 在真旋转、空间反演
和时间反演下的变换规则。

\subsection{投影算符与量子几何张量}

设一组孤立的 \(N\) 条能带（例如费米能级以下的所有占据带），定义投影算符
\begin{equation}
  P(\mathbf{k}) = \sum_{n=1}^{N} \lvert u_{n\mathbf{k}}\rangle\langle u_{n\mathbf{k}}\rvert,
\end{equation}
满足 \(P^2 = P\)，\(P^\dagger = P\)。\textit{量子几何张量}定义为
\begin{equation}
  \mathcal{G}_{ab}(\mathbf{k}) = \operatorname{tr}\bigl[ P(\mathbf{k})\,
  \partial_a P(\mathbf{k})\, \partial_b P(\mathbf{k}) \bigr],
  \label{eq:Gdef}
\end{equation}
其中 \(\partial_a \equiv \partial/\partial k_a\)，迹 \(\operatorname{tr}\) 对能带指标求
迹。其实部和反对称虚部分别给出\textit{量子度规}和\textit{贝里曲率张量}：
\begin{equation}
  \mathcal{G}_{ab} = g_{ab} - \frac{i}{2}\Omega_{ab},
  \qquad
  g_{ab} = \operatorname{Re}\mathcal{G}_{ab},\quad
  \Omega_{ab} = -2\operatorname{Im}\mathcal{G}_{ab}.
\end{equation}
单带 \(n\) 的带归一化量子度规为
\begin{equation}
  G_n^{ab} = \sum_{m\neq n}
  \frac{A_{nm}^a A_{mn}^b + A_{nm}^b A_{mn}^a}{2(\varepsilon_n-\varepsilon_m)},
\end{equation}
其物理变换性质与 \(g^{ab}\) 一致，本文直接用 \(\mathcal{G}_{ab}\) 的变换替代
讨论，不再区分单带与多带。

\subsection{对称操作的分类}

晶体中的对称操作分为两类：
\begin{itemize}
  \item \textbf{幺正操作} \(g\)：如真旋转 \(R\)、空间反演 \(P\)。
  \item \textbf{反幺正操作} \(gT\)：包含时间反演的组合，对波函数附加复共轭作用。
\end{itemize}
任何此二类操作均将动量映射为 \(\mathbf{k} \mapsto R_g\mathbf{k}\)，其中 \(R_g\)
为该操作在动量空间的 \(D\times D\) 正交表示矩阵。投影算符的变换规则为
\begin{align}
  \text{幺正操作 } g &: \; P(R_g\mathbf{k}) = U_g P(\mathbf{k}) U_g^\dagger,
    \label{eq:Pg}\\
  \text{反幺正操作 } gT &: \; P(R_g\mathbf{k}) = U_g P^*(\mathbf{k}) U_g^\dagger
    = U_g P^{\mathsf{T}}(\mathbf{k}) U_g^\dagger,
    \label{eq:PgT}
\end{align}
其中 \(P^*\) 表示矩阵元的复共轭，对厄米算符等于转置。

\subsection{动量导数 \(\partial_a\) 的变换}

引入新坐标 \(\mathbf{k}' = R_g\mathbf{k}\)，由偏导数的链式法则
\[
  \partial'_a
  = \frac{\partial}{\partial k'_a}
  = \sum_{b} \frac{\partial k_b}{\partial k'_a}\,\frac{\partial}{\partial k_b}
  = \sum_{b} (R_g^{-1})^b_{\,a}\,\partial_b
  = \sum_{b} (R_g^{\mathsf{T}})^b_{\,a}\,\partial_b
  = \sum_{b} R_g^{a}{}_b\,\partial_b .
\]
写成算子形式
\begin{equation}
{\partial_{(R_g\mathbf{k})} = R_g\,\partial_{\mathbf{k}} } .
  \label{eq:deriv}
\end{equation}

\textbf{特例}：
\begin{itemize}
  \item 真旋转 \(R\)：\(\partial(R\mathbf{k}) = R\,\partial(\mathbf{k})\)。
  \item 空间反演 \(P\)：\(R_g = -I\)，故 \(\partial_{-k} = -\partial_k\)。
  \item 时间反演 \(T\)：动量同样反转 \(\mathbf{k}\to -\mathbf{k}\)，亦满足
        \(\partial_{-k} = -\partial_k\)。
\end{itemize}

\subsection{量子几何张量的幺正变换}

先考虑幺正操作 \(g\)，其投影算符按式\eqref{eq:Pg}变换。代入定义\eqref{eq:Gdef}：
\begin{align}
  \mathcal{G}_{ab}(R_g\mathbf{k})
  &= \operatorname{tr}\bigl[ P(R_g\mathbf{k})\,
     \partial_{(R_g\mathbf{k})_a} P(R_g\mathbf{k})\,
     \partial_{(R_g\mathbf{k})_b} P(R_g\mathbf{k}) \bigr] \nonumber\\
  &= \operatorname{tr}\bigl[ U_g P U_g^\dagger\,
     \partial_{(R_g\mathbf{k})_a}(U_g P U_g^\dagger)\,
     \partial_{(R_g\mathbf{k})_b}(U_g P U_g^\dagger) \bigr].
\end{align}
利用式\eqref{eq:deriv}替换导数：
\[
  \partial_{(R_g\mathbf{k})_a}(U_g P U_g^\dagger)
  = U_g \bigl( R_g^{a}{}_{a'}\,\partial_{a'}P \bigr) U_g^\dagger
  + \text{含}(\partial U_g)\text{的项}.
\]
代入上一行，利用迹的循环性和 \(U_g^\dagger U_g = I\)，所有含 \(U_g\) 导数的项
在投影算符的幂等性和迹的循环下彼此抵消或投影为零。因此只剩旋转因子的贡献：
\begin{align}
  \mathcal{G}_{ab}(R_g\mathbf{k})
  &= R_g^{a}{}_{a'} R_g^{b}{}_{b'}\,
     \operatorname{tr}\bigl[ P\,\partial_{a'}P\,\partial_{b'}P \bigr] \nonumber\\
  &= R_g^{a}{}_{a'} R_g^{b}{}_{b'}\, \mathcal{G}_{a'b'}(\mathbf{k}).
  \label{eq:Gg_unitary}
\end{align}
这表明 \(\mathcal{G}_{ab}\) 在幺正操作下按二阶张量变换。

\subsection{量子几何张量的反幺正变换}

对反幺正操作 \(gT\)，投影算符按式\eqref{eq:PgT}变换，出现复共轭：
\[
  P(R_g\mathbf{k}) = U_g P^*(\mathbf{k}) U_g^\dagger
  = U_g P^{\mathsf{T}}(\mathbf{k}) U_g^\dagger .
\]
对其求导，复共轭 \(P^*\) 与导数算符对易（导数只涉及 \(\mathbf{k}\) 坐标，不改变
复共轭性质）。结合式\eqref{eq:deriv}，经过和幺正情形相同的抵消途径，
得到
\begin{align}
  \mathcal{G}_{ab}(R_g\mathbf{k})
  &= R_g^{a}{}_{a'} R_g^{b}{}_{b'}\,
     \operatorname{tr}\bigl[ P^*\,\partial_{a'}P^*\,\partial_{b'}P^* \bigr]
   = R_g^{a}{}_{a'} R_g^{b}{}_{b'}\,
     \operatorname{tr}\bigl[ (P\,\partial_{a'}P\,\partial_{b'}P)^* \bigr] \nonumber\\
  &= R_g^{a}{}_{a'} R_g^{b}{}_{b'}\,
     \bigl( \operatorname{tr}\bigl[ P\,\partial_{a'}P\,\partial_{b'}P \bigr] \bigr)^*
   = R_g^{a}{}_{a'} R_g^{b}{}_{b'}\, \mathcal{G}_{a'b'}^*(\mathbf{k}) .
  \label{eq:Gg_anti}
\end{align}
这里用了迹的性质 \(\operatorname{tr}(A^*) = (\operatorname{tr} A)^*\)。

比较式\eqref{eq:Gg_unitary}和\eqref{eq:Gg_anti}可知，反幺正操作比幺正操作
多一个复共轭作用于 \(\mathcal{G}\) 整体，等价于对虚部添加一个负号。

\subsection{量子度规的变换}

量子度规是 \(\mathcal{G}_{ab}\) 的实部：\(g_{ab} = \operatorname{Re}\mathcal{G}_{ab}\)。
从式\eqref{eq:Gg_unitary}和\eqref{eq:Gg_anti}提取实部，复共轭不改变实部，
故幺正和反幺正两种情形得到\textbf{完全相同的变换规则}：
\begin{equation}
{g_{ab}(R_g\mathbf{k}) = R_g^{a}{}_{a'} R_g^{b}{}_{b'}\,g_{a'b'}(\mathbf{k}) } .
\end{equation}

\textbf{三种基本操作的结论}：
\begin{itemize}
  \item \textbf{真旋转 \(R\)}：\(g_{ab}(R\mathbf{k}) = R_a^{a'}R_b^{b'}g_{a'b'}(\mathbf{k})\)。
  \item \textbf{空间反演 \(P\)}：\(R_g=-I\)，两负得正：
        \(g_{ab}(-\mathbf{k}) = g_{ab}(\mathbf{k})\)（偶函数）。
  \item \textbf{时间反演 \(T\)}：亦为偶函数
        \(g_{ab}(-\mathbf{k}) = g_{ab}(\mathbf{k})\)。
\end{itemize}
带归一化量子度规 \(G^{ab}\) 继承相同变换性质。

\subsection{贝里曲率张量的变换}

贝里曲率张量是 \(\mathcal{G}_{ab}\) 的反对称虚部：
\(\Omega_{ab} = -2\operatorname{Im}\mathcal{G}_{ab}\)。
由式\eqref{eq:Gg_anti}，反幺正操作的复共轭对虚部产生一个负号。
统一写作
\begin{equation}
{\Omega_{ab}(R_g\mathbf{k}) = (-1)^{\chi}\,
    R_g^{a}{}_{a'} R_g^{b}{}_{b'}\,\Omega_{a'b'}(\mathbf{k}) },
  \label{eq:Omega_trans}
\end{equation}
其中 \(\chi=0\) 对应幺正操作，\(\chi=1\) 对应反幺正操作。

\textbf{三种基本操作的结论}：
\begin{itemize}
  \item \textbf{真旋转 \(R\)}（幺正）：\(\Omega_{ab}(R\mathbf{k}) = R_a^{a'}R_b^{b'}\Omega_{a'b'}(\mathbf{k})\)。
  \item \textbf{空间反演 \(P\)}（幺正，\(R_g=-I\)）：
        \(\Omega_{ab}(-\mathbf{k}) = \Omega_{ab}(\mathbf{k})\)（偶函数）。
  \item \textbf{时间反演 \(T\)}（反幺正，\(\chi=1\)，\(R_g=I\)）：
        \(\Omega_{ab}(-\mathbf{k}) = -\Omega_{ab}(\mathbf{k})\)（奇函数）。
\end{itemize}

\subsection{贝里曲率赝矢量的变换}

在三维（及二维向 \(z\) 方向推广）中，常将反对称张量 \(\Omega_{ab}\) 对偶为赝矢量
\begin{equation}
  \Omega^c = \frac12\, \epsilon^{cab}\,\Omega_{ab},
  \qquad
  \Omega_{ab} = \epsilon_{abc}\,\Omega^c,
\end{equation}
其中 \(\epsilon^{abc}\) 为 Levi–Civita 符号，\(\epsilon^{xyz}=+1\)。

利用 \(\epsilon\) 在旋转下的恒等式
\[
  \epsilon^{cab} R_g^{a}{}_{a'} R_g^{b}{}_{b'}
  = \det(R_g)\, R_g^{c}{}_{c'}\,\epsilon^{c'a'b'},
\]
对式\eqref{eq:Omega_trans}两边乘 \(\frac12\epsilon^{cab}\) 并对 \(a,b\) 求和，
得到赝矢量的变换规则：
\begin{align}
  \Omega^c(R_g\mathbf{k})
  &= \frac12 \epsilon^{cab}\, (-1)^\chi\, R_g^{a}{}_{a'} R_g^{b}{}_{b'}\,
     \Omega_{a'b'}(\mathbf{k}) \nonumber\\
  &= (-1)^\chi \det(R_g)\, R_g^{c}{}_{c'}\,
     \frac12\epsilon^{c'a'b'}\,\Omega_{a'b'}(\mathbf{k}) \nonumber\\
  &= (-1)^\chi \det(R_g)\, R_g^{c}{}_{c'}\,\Omega^{c'}(\mathbf{k}) .
  \label{eq:Omega_vec}
\end{align}
此即赝矢量版本的通用变换公式。

\textbf{三种基本操作的结论}：
\begin{itemize}
  \item \textbf{真旋转 \(R\)}（\(\chi=0\)，\(\det=+1\)）：
        \(\Omega^c(R\mathbf{k}) = R^{c}_{c'}\,\Omega^{c'}(\mathbf{k})\)。
  \item \textbf{空间反演 \(P\)}（\(\chi=0\)，\(\det=-1\)，\(R_g=-I\)）：
        \(\Omega^c(-\mathbf{k}) = \Omega^c(\mathbf{k})\)（偶函数）。
  \item \textbf{时间反演 \(T\)}（\(\chi=1\)，\(\det=+1\)，\(R_g=I\)）：
        \(\Omega^c(-\mathbf{k}) = -\Omega^c(\mathbf{k})\)（奇函数）。
\end{itemize}

\subsection{汇总表}

表~\ref{tab:sym_trans} 列出动量导数、量子度规、贝里曲率张量与赝矢量在三种
基本操作下的变换规则。

\begin{table}[htbp]
  \centering
  \caption{基本对称操作下各量的完整变换规则}
  \label{tab:sym_trans}
  \footnotesize
  \begin{tabularx}{\textwidth}{c|X|X|X}
    \toprule
    操作 & \(\partial_{\mathbf{k}}\) 的变换 & \(G^{ab}(\mathbf{k})\) 的变换 & \(\Omega^{ab}(\mathbf{k})\) 及 \(\Omega^c(\mathbf{k})\) 的变换 \\
    \midrule
    真旋转 \(R\) &
    \(\partial'_{a} = \sum_{a'} R_{a}^{a'} \partial_{a'} \big|_{R\mathbf{k}}\)
    &
    \(G'^{ab}(R\mathbf{k}) = \sum_{a'b'} R_{a'}^{a} R_{b'}^{b} G^{a'b'}(\mathbf{k})\)
    &
    \(\begin{aligned}
      \Omega'^{ab}(R\mathbf{k}) &= \sum_{a'b'} R_{a'}^{a} R_{b'}^{b} \Omega^{a'b'}(\mathbf{k}),\\
      \Omega'^{c}(R\mathbf{k}) &= \sum_{c'} R_{c'}^{c} \Omega^{c'}(\mathbf{k})
    \end{aligned}\) \\
    \midrule
    空间反演 \(P\) &
    \(\partial'_{a} = -\partial_{a} \big|_{-\mathbf{k}}\)
    &
    \(G'^{ab}(-\mathbf{k}) = G^{ab}(\mathbf{k})\)
    &
    \(\Omega'^{ab}(-\mathbf{k}) = \Omega^{ab}(\mathbf{k})\),\newline
    \(\Omega'^{c}(-\mathbf{k}) = \Omega^{c}(\mathbf{k})\) \\
    \midrule
    时间反演 \(T\) &
    \(\partial'_{a} = -\partial_{a} \big|_{-\mathbf{k}}\)
    &
    \(G'^{ab}(-\mathbf{k}) = G^{ab}(\mathbf{k})\)
    &
    \(\Omega'^{ab}(-\mathbf{k}) = -\Omega^{ab}(\mathbf{k})\),\newline
    \(\Omega'^{c}(-\mathbf{k}) = -\Omega^{c}(\mathbf{k})\) \\
    \midrule
    \(PT\) (组合) &
    \(\partial'_{a} = \partial_{a} \big|_{\mathbf{k}}\)
    &
    \(G'^{ab}(\mathbf{k}) = G^{ab}(\mathbf{k})\)
    &
    \(\Omega'^{ab}(\mathbf{k}) = -\Omega^{ab}(\mathbf{k})\),\newline
    \(\Omega'^{c}(\mathbf{k}) = -\Omega^{c}(\mathbf{k})\) \\
    \bottomrule
  \end{tabularx}
\end{table}
此表是后续非线性电导率在特定磁群下对称性筛选的直接基础。具体的 CrSb 应用实例（非零张量分量、旋转角度依赖的非线性响应）见第~\ref{sec:CrSb_application} 节。
